\documentclass[11pt,a4paper,openany]{book}

% Stile personalizzato Piazza dei Mestieri
\usepackage{stile}

\begin{document}

% ==============================================================================
% COPERTINA MODERNA ED ELEGANTE
% ==============================================================================
\begin{titlepage}
  \newgeometry{top=3cm,bottom=3cm,left=3cm,right=3cm}
  \pagecolor{PiazzaPaper}

  \begin{flushleft}
    {\fontsize{13}{16}\selectfont\bfseries\color{PiazzaPrimary}PIAZZA DEI MESTIERI -- TORINO}\\[0.25em]
    {\fontsize{10}{13}\selectfont\color{PiazzaGray}IeFP -- Operatore Informatico | Prima Annualità (Classe B22-19-20270)}
  \end{flushleft}

  \vspace{3.5cm}

  \noindent{\color{PiazzaPrimary}\rule{\textwidth}{3pt}}
  \vspace{1.5cm}

  \begin{flushleft}
    {\fontsize{34}{40}\selectfont\bfseries\color{PiazzaDark}PROGRAMMAZIONE 1}\\[0.5em]
    {\fontsize{18}{22}\selectfont\color{PiazzaGray}Fondamenti di Logica, Algoritmi e Sviluppo Software}\\[1.5em]
    {\fontsize{11}{14}\selectfont\color{PiazzaPrimary}\bfseries Manuale Didattico Ufficiale del Corso di Laboratorio (80 ore)}
  \end{flushleft}

  \vspace{1.5cm}
  \noindent{\color{PiazzaPrimary!40}\rule{\textwidth}{1pt}}

  \vfill

  \begin{flushleft}
    \textbf{\color{PiazzaDark}Docente:}\hspace{1.5em}{\color{PiazzaDark}Daniele Cerrina}\\
    \textbf{\color{PiazzaDark}Email:}\hspace{2.4em}{\small\texttt{daniele.cerrina@immaginazioneelavoro.it}}\\
    \textbf{\color{PiazzaDark}Anno Formativo:}\hspace{0.5em}2026 / 2027
  \end{flushleft}

  \restoregeometry
\end{titlepage}
\pagecolor{white}

% ==============================================================================
% PARTE INTRODUTTIVA (FRONTMATTER)
% ==============================================================================
\frontmatter

% Pagina del Colophon / Info sul corso
\thispagestyle{empty}
\vspace*{12cm}
\begin{flushleft}
  \small\color{PiazzaGray}
  \textbf{Programmazione 1 -- Fondamenti e Laboratorio}\\
  Percorso triennale IeFP per la Qualifica di Operatore Informatico.\\
  \textcopyright~2026--2027 Piazza dei Mestieri / Immaginazione e Lavoro.\\
  Docente: Daniele Cerrina.\\
  \vspace{0.8em}
  Documento didattico in continuo aggiornamento durante l'anno formativo.\\
  Compilato con \LaTeX\ e \textsf{pdflatex}.
\end{flushleft}

\clearpage

% Introduzione e Presentazione del Corso
\chapter*{Presentazione del Libro}
\addcontentsline{toc}{chapter}{Presentazione del Libro}

Questo volume rappresenta il punto di riferimento e il libro di testo per il corso di \textbf{Programmazione 1} (80 ore di laboratorio) presso la Piazza dei Mestieri di Torino.

L'approccio adottato è fortemente pratico e incrementale: ogni lezione svolta in aula e in laboratorio viene consolidata e integrata all'interno di questo testo, fornendo agli allievi:
\begin{itemize}
  \item Le spiegazioni teoriche chiare dei concetti chiave.
  \item Gli esempi concreti e analogie vicine all'esperienza quotidiana e videoludica.
  \item Gli schemi logici ed esercizi svolti e da svolgere.
  \item Il percorso di sviluppo del \textbf{gioco guida annuale} votato dalla classe.
\end{itemize}

\vspace{1em}
\noindent\textbf{Come consultare questo documento:}
Il documento è provvisto di un \textbf{indice ipertestuale e cliccabile}. Cliccando su qualsiasi capitolo o sezione nell'Indice Generale o nei segnalibri del lettore PDF, verrai indirizzato direttamente alla pagina corrispondente.

\clearpage

% INDICE GENERALE CLICCABILE
\pdfbookmark[0]{\contentsname}{toc}
\tableofcontents
\clearpage

% ==============================================================================
% PARTE PRINCIPALE (MAINMATTER)
% ==============================================================================
\mainmatter

% ------------------------------------------------------------------------------
% MODULO 1: IL PENSIERO COMPUTAZIONALE E GLI ALGORITMI
% ------------------------------------------------------------------------------
\part{Pensiero Computazionale e Algoritmi}

% Capitolo 1: Prima lezione del 21 Settembre 2026
\chapter{Introduzione agli Algoritmi e Pensiero Computazionale}
\label{chap:intro-algoritmi}

\begin{quote}
\itshape
\small
``Un algoritmo deve essere visto per essere creduto, ma soprattutto per essere compreso: è il ponte invisibile che trasforma un problema disordinato in una soluzione esatta.''
\end{quote}
\vspace{1em}

\noindent\textbf{Riferimento Lezione:} 21 Settembre 2026 -- Prima Lezione di Laboratorio\\
\textbf{Competenze chiave:} Comprensione del modello Input/Output, identificazione delle proprietà di un algoritmo, formulazione di istruzioni non ambigue.

\vspace{1.5em}

% ------------------------------------------------------------------------------
\section{Che cos'è un Algoritmo?}
% ------------------------------------------------------------------------------

Nel linguaggio comune siamo abituati a risolvere problemi ogni giorno senza pensare alle singole azioni che compiamo. 
Quando però vogliamo affidare la risoluzione di un problema a un computer (o spiegare a un'altra persona come agire senza margini d'errore), abbiamo bisogno di un metodo rigoroso.

\begin{concetto}[Definizione di Algoritmo]
Un \textbf{algoritmo} è una sequenza finita e ordinata di istruzioni elementari, chiare e non ambigue, la cui esecuzione porta alla risoluzione di un determinato problema in un tempo finito (altrimenti parleremo di loop, argomento di cui ci occuperemo più avanti).
\end{concetto}

In informatica ogni problema viene formalizzato attraverso il modello fondamentale di elaborazione:

\begin{itemize}
  \item \textbf{Input (Dati di ingresso):} Rappresenta l'informazione o la condizione iniziale fornita all'algoritmo. È ciò che ``entra'' nel processo.
  \item \textbf{Elaborazione:} L'insieme ordinato di istruzioni elementari applicate ai dati di ingresso.
  \item \textbf{Output (Risultato):} Il dato finale o l'azione risultante al termine dell'elaborazione. È ciò che ``esce'' dal processo.
\end{itemize}

\vspace{1em}
\begin{center}
\begin{tcolorbox}[colback=white,colframe=PiazzaGray!40,width=0.88\textwidth,arc=2mm,halign=center]
\textbf{INPUT} (Dati / Problema) $\longrightarrow$ \colorbox{PiazzaPrimary!10}{\textbf{~[ ALGORITMO ]~}} $\longrightarrow$ \textbf{OUTPUT} (Risultato / Soluzione)
\end{tcolorbox}
\end{center}

% ------------------------------------------------------------------------------
\subsection{Esempi dal Quotidiano e dal Mondo Digitale}
% ------------------------------------------------------------------------------

Gli algoritmi non appartengono solo al codice sorgente: li sperimentiamo quotidianamente in molteplici contesti:

\begin{esempio}[Algoritmi nella vita reale]
\begin{enumerate}
  \item \textbf{Tutorial per costruire una casa in Minecraft:} Un elenco ordinato di azioni (1. Raccogli 64 blocchi di legno, 2. Costruisci una base $6 \times 6$, 3. Innalza i muri per 4 blocchi, 4. Posiziona la porta e le torce). Se inverti i passaggi (ad esempio posizionare il tetto prima dei muri), l'obiettivo fallisce.
  \item \textbf{Ricetta di cucina (es. Torta al cioccolato):} L'elenco degli ingredienti rappresenta l'\textbf{Input}; la sequenza dei passaggi (pesare, mescolare, infornare a 180 gradi per 35 minuti) costituisce l'\textbf{Algoritmo}; il dolce pronto è l'\textbf{Output}.
  \item \textbf{Istruzioni di montaggio (LEGO o mobili IKEA):} Diagrammi e numeri sequenziali guidano la costruzione partendo dai singoli pezzi fino all'oggetto completo.
\end{enumerate}
\end{esempio}

% ------------------------------------------------------------------------------
\section{Le Quattro Proprietà Fondamentali}
% ------------------------------------------------------------------------------

Non ogni elenco di frasi o istruzioni può essere definito un algoritmo valido. Per essere considerato formalmente \textbf{corretto}, un algoritmo deve soddisfare quattro caratteristiche essenziali:

\begin{enumerate}
  \item \textbf{Finitezza:} L'algoritmo deve necessariamente terminare dopo aver eseguito un numero finito di istruzioni. Un procedimento che gira all'infinito senza mai giungere a una conclusione non è un algoritmo.
  
  \item \textbf{Definitezza (Non ambiguità):} Ogni istruzione deve essere univoca e comprensibile senza alcuna interpretazione soggettiva. Frasi come ``mescola un po'\,'' o ``aspetta qualche secondo'' non sono adatte a un algoritmo informatico: bisogna specificare ``mescola per 3 minuti a velocità 2'' o ``attendi 5 secondi''.
  
  \item \textbf{Efficienza:} L'algoritmo deve risolvere il compito impiegando una quantità *ragionevole* di tempo e di memoria. 
  \begin{consiglio}[Nota]
  Nel corso di quest'anno la priorità assoluta sarà comprendere la \textbf{logica} e raggiungere il \textbf{corretto funzionamento}; la misurazione scientifica dell'efficienza (complessità computazionale) sarà affrontata negli anni successivi.
  \end{consiglio}

  \item \textbf{Generalità:} L'algoritmo non deve risolvere un singolo caso specifico e isolato, ma deve essere in grado di risolvere \textbf{qualsiasi problema appartenente a quella specifica classe}.
\end{enumerate}

\begin{concetto}[Il Principio della Generalità]
\textbf{Il nostro obiettivo annuale è imparare a risolvere qualsiasi problema della categoria che stiamo affrontando.}

\vspace{0.5em}
\textbf{Esempio pratico:} Se progettiamo un algoritmo per calcolare la media aritmetica di 7 numeri, l'algoritmo deve funzionare correttamente sia con la sequenza $[1, 2, 3, 4, 5, 6, 7]$, sia con $[9, 8, 7, 6, 5, 4, 3]$, sia con qualsiasi altra combinazione di 7 valori numerici, inclusi numeri decimali o negativi.
\end{concetto}

% ------------------------------------------------------------------------------
\newpage
\section{Esercizio Guidato: Attraversamento Pedonale}
% ------------------------------------------------------------------------------

Proviamo a scrivere un algoritmo esatto per un problema che affrontiamo ogni giorno.

\begin{esercizio}[Attraversare una strada con un semaforo]
Scrivere l'algoritmo passo-passo, il più rigoroso e aderente possibile alle quattro proprietà esaminate, per risolvere il seguente problema:\\
\textbf{Problema:} Mi trovo sul marciapiede davanti alle strisce pedonali e devo attraversare la strada in presenza di un semaforo pedonale. Cosa faccio?
\end{esercizio}

\begin{soluzione}[Algoritmo con Scelta Condizionale e Ripetizione]
Ecco la formalizzazione dell'algoritmo divisa in istruzioni elementari:

\begin{enumerate}
  \item \textbf{Osservazione:} Guardo il semaforo pedonale.
  \item \textbf{Controllo condizionale:}
  \begin{itemize}
    \item \textbf{SE} il semaforo è \textbf{VERDE}:
    \begin{enumerate}
      \item Inizio ad attraversare a passo costante sulle strisce pedonali.
      \item Faccio un passo dopo l'altro \textbf{FINCHÉ} non ho raggiunto il marciapiede opposto.
      \item Fine dell'algoritmo (Obiettivo raggiunto con successo!).
    \end{enumerate}
    \item \textbf{ALTRIMENTI} (il semaforo è rosso o giallo):
    \begin{enumerate}
      \item Resto fermo sul marciapiede di partenza in sicurezza.
      \item Aspetto che il semaforo cambi stato.
      \item \textbf{Torno al punto 1} (ripeto l'osservazione).
    \end{enumerate}
  \end{itemize}
\end{enumerate}
\end{soluzione}

\begin{attenzione}[Attenzione all'ambiguità e ai cicli infiniti!]
\begin{itemize}
  \item Se dimentichiamo la condizione ``\textbf{FINCHÉ} non arrivo dall'altra parte'', una persona programmata come un computer continuerebbe a camminare dritto all'infinito anche dopo aver superato la strada!
  \item Se omettiamo di dire cosa fare nel caso il semaforo sia rosso, il calcolatore si bloccherebbe per mancanza di istruzioni.
  \item Notiamo come questo algoritmo funzioni con qualsiasi strada che abbia un semaforo in cui c'è almeno il verde.
\end{itemize}
\end{attenzione}




% ------------------------------------------------------------------------------
% SEZIONI FUTURE (Man mano che il corso prosegue, scommenta o aggiungi qui i file)
% ------------------------------------------------------------------------------
% \input{capitoli/02_diagrammi_di_flusso}
% \input{capitoli/03_variabili_e_tipi}

% \part{Costrutti di Controllo e Flusso di Esecuzione}
% \input{capitoli/04_strutture_condizionali}
% \input{capitoli/05_cicli_e_iterazioni}

% \part{Modularità e Strutture Dati}
% \input{capitoli/06_funzioni_e_procedure}
% \input{capitoli/07_vettori_e_collezioni}

% \part{Il Gioco Guida Annuale}
% \input{capitoli/08_progetto_videogioco}

% ==============================================================================
% PARTE CONCLUSIVA (BACKMATTER)
% ==============================================================================
\backmatter

\chapter{Guida Rapida alla Sintassi}
\noindent In questa appendice finale verranno raccolti durante l'anno i promemoria di sintassi, le parole chiave e i comandi più utilizzati in laboratorio.

\end{document}
